\chapter{Introduction}
\label{Chapter1}

Embedded systems are widely deployed in safety- and security-critical settings such as
industrial control, automotive platforms, and consumer devices. In these environments,
adversaries may obtain physical access and attempt unauthorized firmware replacement or
invasive debugging. Secure boot mechanisms authenticate the initial firmware image, but
they do not necessarily restrict \emph{post-boot} programming and debug access. As a
result, debug/test interfaces remain a high-impact attack surface because they expose
memory and execution control and can bypass software-level protections
(\cite{ieee1149_1,ali2014boundaryscanattack}). Recent work has proposed authenticated
debug mechanisms, yet such controls often remain external to the execution model
(\cite{lapeyre2022jtagauth}).

This thesis addresses the gap by enforcing device-local authorization before execution
and debug enablement. The proposed approach introduces a minimal ISA extension and a
custom instruction (\texttt{kc}) that triggers verification against key material stored
in protected non-volatile memory. An on-device authenticator FSM gates both normal
execution and debug access, ensuring that privileged capabilities are released only
after successful authorization. The approach is integrated into the toolchain so that
authorization is emitted and preserved through compilation, enabling reproducible and
auditable binaries. The design targets practical deployment by leveraging ISA
extensibility (e.g., RISC-V) and maintaining compatibility with standard build flows
(\cite{waterman2014riscv,bin2020bootrom,rv2018securebootreview}).

\section{Motivation}
Physical access threats remain common for embedded devices deployed in the field. When
debug interfaces are exposed or weakly protected, adversaries can extract secrets,
modify firmware, or bypass runtime policies. Conventional secure boot verifies an image
at startup but does not necessarily prevent post-boot debug access or reprogramming
events. This motivates a mechanism that binds authorization to the execution model
itself, rather than relying solely on host-side policy or external tooling. A hardware-
software co-design with an ISA-visible authorization primitive provides a practical
path to enforce such a policy directly on the device.

\section{Problem Statement}
Current embedded defenses leave three gaps. First, post-boot debug and programming
access can be enabled without a device-local authorization check. Second, key material
used for authorization is often accessible through ordinary memory transactions or
insecure provisioning flows, enabling extraction by untrusted firmware. Third, host-
side policies are insufficient when an attacker can replace firmware or bypass the
host toolchain. A device-resident authorization mechanism is required that: (i) gates
execution and debug access on-device, (ii) keeps keys out of the normal load/store
path, and (iii) is integrated into the compiler/toolchain to ensure enforceability and
auditability.

\section{Scope and Objectives}
This thesis focuses on an embedded platform with a minimal ISA extension and a
hardware verification datapath. The scope includes: (i) definition of the custom
instruction \texttt{kc}, (ii) protected key storage in non-volatile memory, (iii) an
authenticator FSM that controls execution and debug gating, and (iv) toolchain support
to emit and preserve the authorization primitive during compilation. The scope does
not include full cryptographic key management at scale, advanced side-channel defenses,
or production-grade certification workflows.

The primary objectives are:
\begin{itemize}
  \item Design an ISA-level authorization primitive and its hardware datapath.
  \item Enforce device-local gating of execution and debug access with fail-stop
        behavior on authorization failure.
  \item Integrate the authorization instruction into the compiler/toolchain pipeline.
  \item Validate end-to-end functionality using build artifacts and host workflow logs.
\end{itemize}

\section{Contributions}
The key contributions of this thesis are:
\begin{itemize}
  \item A minimal ISA extension that triggers device-local authorization via
        \texttt{kc} without exposing key material to the normal memory hierarchy.
  \item A security FSM that gates both execution and debug enablement and enforces
        fail-stop behavior on failed authorization.
  \item A compiler/toolchain integration that emits and preserves the authorization
        primitive, enabling reproducible and inspectable binaries.
  \item An end-to-end validation workflow demonstrating authorized programming and
        debug enablement only after successful key verification.
\end{itemize}

\section{Thesis Structure}
The remainder of the thesis is organized as follows:
\begin{itemize}
  \item \textbf{Chapter 2} reviews related work in secure boot, debug interface
        security, trusted execution, and toolchain support.
  \item \textbf{Chapter 3} presents the system architecture, including the ISA
        extension, protected key storage, and authenticator FSM.
  \item \textbf{Chapter 4} describes the toolchain integration and host workflow.
  \item \textbf{Chapter 5} reports evaluation results and evidence artifacts.
  \item \textbf{Chapter 6} discusses implications, limitations, and future directions.
\end{itemize}
